%% kapa: exam=1 round=37 year=2026 period=2 form=A subject=accounting subject_ko=회계학 seq=22 number=62
%% source: Q-Net 감정평가사 기출문제 (cst00309) artlSeq=5258924 file=raw/1008/5258924/2026년 제37회 감정평가사 1차 2교시.pdf
%% conversion: kordoc 4.24.51 -> tools/gen1.py (auto); review=proofread-sonnet; flags=none
\begin{question}{62}
(주)감평은 20×1년 초에 부여일로부터 3년의 지속적인 용역제공을 조건으로 직원 100명에게 주식선택권을 1인당 10개씩 부여하였다. 20×1년 초 주식선택권의 단위당 공정가치는 ￦100이며, 주식선택권은 20×4년 초부터 행사할 수 있다. (주)감평의 연도 말 실제 누적 퇴직자 수 및 추가퇴직 예상자 수는 다음과 같다.
\par
{\small\centering
\begin{tabular}{|c|c|c|}\hline
구분 & 실제 누적 퇴직자 수 & 추가퇴직 예상자 수 \\\hline
20×1년 말 & 3명 & 13명 \\\hline
20×2년 말 & 7명 & 12명 \\\hline
\end{tabular}\par}
(주)감평은 20×2년 초에 주식선택권의 행사가격을 낮추어 종업원에게 유리한 조건변경을 하였으며, 이러한 조건변경으로 주식선택권의 단위당 공정가치가 ￦20 증가하였다. 20×2년도에 인식할 주식보상비용은? \par
\begin{choices}
  \item ￦26,000
  \item ￦28,000
  \item ￦32,400
  \item ￦34,100
  \item ￦36,800
\end{choices}
\end{question}
\answer{4}
